\section{Selected ordinary multiplicities in rank three}
\label{sec:rank-three-multiplicity}
\hypertarget{source-ordinary-rank-three-multiplicity}{}

Existence of a classical eigensystem sufficed to start the rank-four comparison. In the rank-three trace calculation the coefficient of each contributing eigensystem is its multiplicity in the specified ordinary module. Proposition~\ref{prop:ordinary-three-local-dimensions} computes the three local factors: the ordinary multiplicity above \(\ell\), the multiplicity of the prescribed tame type, and the contribution of the depth-one projector \eqref{eq:depth-one-projector-splitting} at the exceptional place.

Let \(L/F\) be the CM extension in the rank-three ordinary
setting, unramified at every finite place of \(F\), and let
\(\ell\) be the coefficient prime. Use the definite unitary group of rank three for \(L/F\), compact at each real place, with the standard maximal compact subgroups at the finite places. Keep the
chosen split places for the level and their orientations
fixed. Denote the exceptional place of \(F\) by \(v_0\) and the chosen place of \(L\) above it by \(w\). Splitting and the residue-characteristic hypothesis require
\[
 L\otimes_F F_{v_0}\simeq L_w\times L_{\bar w},\qquad L_w\simeq F_{v_0},\quad \operatorname{char} k(v_0)\ne\ell .
\] Its level is
principal congruence of depth one. At the other split tame
places the coefficient type is fixed. The places above
\(\ell\) in the ordinary charts are split.

Fix an arithmetic weight \(\lambda\), a level and coefficient
type in this setting, and an eigensystem \(x\) of the
localized ordinary Hecke module. Assume that its arithmetic
eigenspace is nonzero and that its specialized determinant
is the determinant of an absolutely irreducible
representation. The localized ordinary operators have
unit eigenvalues. We compare the multiplicity of \(x\)
with that of another such point at the same algebraic weight
exponents and level. Taking both multiplicities on the image of the same integral idempotent makes the comparison depend on dimensions at fixed level and weight.

\subsection{The stable count at a fixed weight}
\label{sec:rank-three-stable-count}

The multiplicity has a global factor determined by the archimedean weight. At a cuspidal base change, the definite group has a unique contributing representation of the prescribed weight that is spherical at the inert places. Its multiplicity can therefore be determined before imposing the ordinary and tame projectors.

\begin{proposition}[The stable count at the fixed algebraic weight]
\label{prop:ordinary-stable-constituent-count}
For rank \(1\le n\le4\), let
\(\lambda^\infty\) be a fixed dominant weight at the
infinite places. Let \(\pi\) be a regular algebraic
conjugate self-dual cuspidal representation of
\(\mathrm{GL}_n(\mathbb A_L)\), with the infinitesimal
character prescribed by \(\lambda^\infty\), spherical at
the nonsplit finite places.

At a factorizable finite
level \(K\) that is standard maximal compact at nonsplit
places, write \(\sigma_i\), \(i\in I\), for the
irreducible unitary constituents with the prescribed
weight in the finite-level spectrum. The index set is
finite and counts the summands of a direct sum of
level-fixed spaces. Define \(I_\pi\) by agreement of the
spherical Satake multisets of \(\sigma_i\) and \(\pi\) at
all but finitely many split places. Assume the split
components of \(\pi\) have vectors fixed by the
prescribed split level. Then
\[
                   |I_\pi|=a_0(\lambda^\infty)>0,
\]
where \(a_0\) is chosen for \(L,n,\lambda^\infty\)
before choosing \(\pi\), the places of the level, or
the coefficient system.

\end{proposition}

\begin{proof}
\proofstep{Identify the base change from the split Hecke operators.}
Let \(F=L^+\) and let \(\sigma_i\) be one of the summands in \(I_\pi\). Since the group is definite, its real components belong to the discrete series. It is quasi-split at the finite places in the standard self-dual lattice model: at an inert place the extension is unramified and the standard Hermitian lattice is unimodular. Its stabilizer is hyperspecial. In particular the fixed vectors used at a nonsplit place belong to
\[
                   \sigma_{i,v}^{G(\mathcal O_{F_v})}.
\] The field \(F\) has degree at least two since an imaginary quadratic extension of \(\mathbb Q\) cannot be unramified at every finite prime.

Base change of \(\sigma_i\) gives a sum \(\Pi_i\) of discrete general linear constituents. For a discrete constituent of cuspidal rank \(m\) and segment length \(s\), the twists of its \(s\) summands in the rank-\(n\) Galois normalization are
\[
                         (m-n+s-1)/2-a,\qquad 0\le a<s,
\]
as in Section~\ref{sec:count-two-level-attachments}. Take their direct sum over all constituents. At all but finitely many places of \(L\) these summands recover the local unramified parameter of \(\Pi_i\). The algebraic regularity of the twists and their polarizations follow from the fixed algebraic infinitesimal character as in that section. \cite[Corollary~5.3]{labesse-cm} \cite[Theorem~2.1.1]{blggt}

Compare with the similarly normalized attachment to \(\pi\). Agreement at the chosen half of a split place gives agreement at its conjugate by conjugate self-duality, with the same rank normalization on both sides. Hence the Frobenius polynomials agree at almost all places \(w\) split over \(F\). At an inert place above \(v\) the norm is \((\mathrm N v)^2\), so the omitted set has Dirichlet density zero:
\[
 \sum_{w\ {\rm inert}}(\mathrm Nw)^{-z}\le\sum_v(\mathrm Nv)^{-2z}=O(1)\qquad(z\downarrow1).
\] Chebotarev density and Brauer--Nesbitt therefore give
\[
          r(\Pi_i)^{\rm ss}\simeq r(\pi)^{\rm ss}.
\]
Reading the characteristic polynomials at all their common unramified primes now gives equality of the almost-everywhere local parameters of \(\Pi_i\) and \(\pi\). Uniqueness of the general linear cuspidal support, with the segment exponents included, gives \(\Pi_i=\pi\). This is the general linear rigidity used in the proof of Labesse's multiplicity formula. \cite[proof of Theorem~5.4]{labesse-cm} \cite{moeglin-waldspurger}

\proofstep{Recover the constituent from the cuspidal base change.}
At a split place, base change identifies the component at the marked half with the unitary representation. At an inert place, \(K_v\)-invariance selects the spherical representation. Spherical base change is injective there: in the standard diagonal parametrization the conjugacy class is described by unordered inverse pairs, and, in odd rank, the additional middle entry. The base-change multiset recovers these pairs. At a real place of the compact group the dominant highest weight determines the representation uniquely. Consequently any two indices in \(I_\pi\) satisfy
\[
         \sigma_{i,v}\simeq\sigma_{i',v}\ (\text{all }v)
           \quad\Longrightarrow\quad \sigma_i\simeq\sigma_{i'}.
\]

It remains to see that this representation occurs, with the required hyperspecial vectors. In the refined cuspidal trace identity in the proof of Labesse's Theorem~5.4, take the hyperspecial unit at every inert place. With the twisted intertwiner normalized to fix the spherical vector, its twisted trace on the projected line is one. At a split place transfer exchanges the two factors. On finite-dimensional images the identity
\[
 \operatorname{tr}\big((A\otimes B)\operatorname{swap}\big)=\operatorname{tr}(AB)
\]
follows by summing diagonal coefficients in a basis. Thus the matched level units have twisted trace \(\dim\pi_w^{K_v}>0\). At infinity use the discrete-series test of the prescribed weight. Once the general linear eigensystem is isolated, the twisted cuspidal term is nonzero; the discrete unitary sum therefore contains a constituent with these level and weight conditions. For the fields \(L=L^+K_0\) of \eqref{eq:constituent-cm-field}, this existence, including the hyperspecial condition, is also exactly Thorne's descent theorem. \cite[Theorem~5.4 and its proof]{labesse-cm} \cite[Theorem~4.4]{thorne-unitary-descent}

\proofstep{Distinguish multiplicity from the dimension of the fixed space.}
The unique constituent is cuspidal on the definite group and spherical at the inert places. Compactness at infinity supplies the discrete-series and ellipticity conditions in Labesse's multiplicity-one theorem. Hence its automorphic multiplicity is one. In the finite-level direct sum its contribution is one summand \(\sigma_f^K\), of possibly larger dimension. Thus, for its unique contributing isomorphism class \([\sigma]\),
\[
                |I_\pi|=m_{\rm cusp}(\sigma)=1 .
\] The local dimensions used below are computed inside that summand. \cite[Theorem~5.4]{labesse-cm}
\end{proof}

\subsection{Which constituents contribute to the ordinary system}

The stable identity counts a packet. The ordinary trace only sees those constituents in the packet that realize the chosen eigensystem with the prescribed local characters.

\begin{lemma}[Constituents with the selected ordinary eigensystem]
\label{lem:ordinary-contributing-constituents}
For the nonzero ordinary eigenspace with an absolutely irreducible
specialization, the contributing constituents are precisely the
members of \(I_{\pi_x}\) with the prescribed local ordinary character.
If \(i_*\) indexes one contributing constituent, then
\[
 i\in I_{\pi_x}\quad\Longrightarrow\quad
     \sigma_{i,v}\simeq\sigma_{i_*,v}\qquad(v\text{ split}).
\]
These split components give the ordinary line, tame type and
projected exceptional factor for each such constituent.
\end{lemma}

\begin{proof}
\proofstep{Control at a common classical depth.}
The characteristic-zero specialization at \(\lambda\)
identifies the integral ordinary fiber with a space of
classical forms with a prescribed algebraic coefficient
representation. Let \(\lambda^\infty\) be its weight
under the coefficient-field isomorphism with
\(\mathbb C\). The finite-level decomposition at this
weight may be taken at a common depth that kills the
tame type on the deeper subgroup and contains the
specified ordinary invariants. If \(c'_0\) bounds the tame-type level and \(b\) is the classical specialization's level parameter, the depth chosen for this test is
\[
 c'_0+b+1 .
\] Expansion in the
finite-level direct sum gives the direct sum of
coefficient subspaces \(Y_i\) on the classical
space. Here a constituent is first restricted to the
right level and then paired with the algebraic and
tame coefficient representations. The resulting
subspaces are invariant under the Hecke operators,
the level constraints and the exceptional-place
operator.

Write \(Z_x\) for the joint generalized eigenspace
with the prescribed ordinary conditions. The nonzero eigenspace meets at least one summand. Choose \(i_*\) and a vector with
\[
 0\ne v\in Y_{i_*}\cap Z_x .
\] The characteristic
polynomials on this summand are those of the
specialized determinant, by the unramified Hecke
polynomial identity. Equality on a dense set of
Frobenius elements gives equality everywhere,
and Brauer--Nesbitt identifies the semisimple
representation attached to the summand with the
given irreducible representation. The rank-three cuspidal transfer therefore attaches
to \(\sigma_{i_*}\) a regular algebraic conjugate
self-dual cusp \(\pi_x\) of the prescribed
infinitesimal character. It is unramified at
nonsplit places, where the unitary constituent
has maximal-compact fixed vectors. At each chosen split place \(v\), denote the prescribed group at the common depth by \(K_v\). The decomposition of the constituent at this level gives
\[
 (\sigma_{i_*,v})^{K_v}\ne0 .
\]

If \(Y_i\cap Z_x\ne0\), use the cofinite split
places outside the level and the ordinary set.
At each such place the elementary Hecke generators
act as scalars on the corresponding spherical
line. For each elementary Hecke generator \(T_{w,j}\) let
\(a_{w,j}\) be its eigenvalue in the specified system and
\(b_{w,j}\) the scalar on the constituent. A nonzero
\(v\in Y_i\cap Z_x\) satisfies
\[
 0=(T_{w,j}-a_{w,j})^N v=(b_{w,j}-a_{w,j})^N v
\]
for some \(N\). Thus \(b_{w,j}=a_{w,j}\) for every \(j\).
The Hecke polynomials agree with the same normalization.
Since the Hecke polynomial recovers the multiset
of Satake parameters, this gives \(i\in I_{\pi_x}\).

\proofstep{The converse by the split component comparison.}
One inclusion identifies the packet of every contributing form. For the converse, comparison at the split places must show that packet members with the prescribed local character really contribute to the localized module.

Conversely the split component comparison shows that membership in the packet identifies the local components, namely
\[
 i\in I_{\pi_x}\quad\Longrightarrow\quad \sigma_{i,v}\simeq\sigma_{i_*,v}\qquad(v\text{ split}).
\] In particular the spherical
eigenvalues match every tested spherical operator.
Here the sets of split level primes coincide, \(S'=S\), so no additional level variables occur. Within the classical
ordinary conditions, the only remaining test is
the ordinary joint generalized eigenspace of the
local components at the places above \(\ell\).
Its rank is at most one for the prescribed
normalization and is positive for \(i_*\), by the
ordinary calculation below.

All members of the packet have the same local components at these places. With \(r_i\) denoting the ordinary rank just bounded,
\[
 1\le r_{i_*}\le1,\qquad r_i=r_{i_*}=1\quad(i\in I_{\pi_x}).
\]
The factorization of the classical ordinary
space into local components now identifies
\(Y_i\cap Z_x\), with the exceptional projector
inserted, with the tensor product of that
ordinary line, the tame type space and the
selected exceptional space. Its dimension is
therefore independent of \(i\in I_{\pi_x}\).

The same argument with the identity in place
of the exceptional projector shows that the
tame type space and the exceptional
principal-congruence fixed space both have
positive dimension. Indeed their tensor
product with the ordinary line at \(i_*\)
contains a nonzero vector. This supplies the
nonvanishing hypothesis for projecting the
exceptional space.

\end{proof}

\subsection{The three local dimensions}

For the surviving constituents, the projected dimension is a product of local dimensions. The calculation separates the ordered ordinary line, the prescribed tame type and the exceptional split place, where the depth-one projector changes the dimension.

\begin{proposition}[The ordinary, tame and depth-one dimensions]
\label{prop:ordinary-three-local-dimensions}
For each constituent in Lemma~\ref{lem:ordinary-contributing-constituents}, the ordinary joint
eigenspaces at the places over \(\ell\) are lines. The fixed tame
type has dimension \(d_{\mathrm{tame}}>0\) independent of the point
at fixed tame labels. At the exceptional split place the prescribed
depth-one projector has rank 
\[
 D_0=(q-1)^{n-1}
\]
 on the fixed
vectors, with \(q\) the residue cardinality.
\end{proposition}

\begin{proof}
\proofstep{The depth-one projector at the exceptional place.}
At \(v_0\) write \(K_1\) for principal congruence in
\(K=\mathrm{GL}_n(\mathcal O)\), where \(\mathcal O\)
is the local integer ring. Let \(k\) have order \(q\)
and let \(U\) be the upper unitriangular group. Fix a
nontrivial additive character \(\bar\psi\) of \(k\).
The characters
\[
  \chi_t(u)=\bar\psi\!\left(\sum_{i=1}^{n-1}t_i u_{i,i+1}\right),
  \qquad t\in(k^\times)^{n-1},
\]
are distinct characters of \(U(k)\). Their projectors
sum to the convolution idempotent
\[
 e=\frac{1}{|U(k)|}\sum_{u\in U(k)}
    \prod_{i=1}^{n-1}\bigl(q\,\mathbf{1}_{u_{i,i+1}=0}-1\bigr)[u].
\]
For a simple entry \(x\in k\), additive orthogonality gives
\[
 \sum_{t\in k^\times}\bar\psi(-tx)
 =\sum_{t\in k}\bar\psi(-tx)-1=q\,\mathbf{1}_{x=0}-1.
\]

Write \(e_t=|U(k)|^{-1}\sum_u\chi_t(u)^{-1}[u]\).
Character orthogonality on the finite group gives
\begin{equation}\label{eq:depth-one-projector-splitting}
 e_t e_{t'}=\mathbf 1_{t=t'}e_t,\qquad
 e=\sum_{t\in(k^\times)^{n-1}}e_t,\qquad
 eW^{K_1}=\bigoplus_t W^{K_1}[\chi_t].
\end{equation}
Indeed the sum over the independent \(t_i\) factors into the simple-entry
sums just evaluated. The denominator \(|U(k)|=q^{n(n-1)/2}\) is
a unit in the coefficient ring at \(\ell\). The group acts through compact
lifts on \(K_1\)-fixed vectors. In the ordinary module
these translations are at a place distinct from the
spherical and ordinary operators, so \(e\) commutes
with their actions.

Each contributing component \(W\) at \(v_0\) is a
local component of a cuspidal transfer. In the Kirillov model the translates of a Whittaker functional \(\Lambda\) by the mirabolic group \(P\) separate vectors. Equivalently,
\[
 \bigcap_{g\in P}\ker\bigl(v\mapsto\Lambda(gv)\bigr)=\{0\}.
\]
\cite[\S5.2 and Corollary 8.2]{zelevinsky}
Take the additive character on the local field of
conductor the maximal ideal and use lifts of the
\(t_i\) to units. Diagonal conjugation within the
mirabolic group gives the separating functional
for this choice of \(\chi_t\).

Starting from \(0\ne w\in W^{K_1}\), choose a
translate with nonzero Whittaker evaluation. By
Iwasawa decomposition and equivariance it can be
written \(D w_1\), where \(D=\operatorname{diag}(d_i)\)
and \(w_1\in W^{K_1}\). Write \(x_{ij}(a)=1+aE_{ij}\), let \(\mathfrak p\) be the
maximal ideal and let \(\widetilde t_i\) be the unit lift of \(t_i\).

For \(a\in\mathfrak p\), the simple root group fixes \(w_1\). Thus
\[
 \Lambda(Dw_1)
 =\Lambda(Dx_{i,i+1}(a)w_1)
 =\psi(\widetilde t_i d_i d_{i+1}^{-1}a)\,\Lambda(Dw_1).
\]
The evaluation is nonzero and \(\psi\) has conductor \(\mathfrak p\);
hence \((d_i/d_{i+1})\mathfrak p\subseteq\mathfrak p\).
Consequently \(v(d_1)\ge\cdots\ge v(d_n)\). For \(i>j\), this gives
\[
 D^{-1}x_{ij}(a)D=x_{ij}(a d_j/d_i)\in K_1
       \qquad(a\in\mathfrak p).
\]
So the lower and diagonal parts of \(K_1\) fix \(Dw_1\);
the upper part preserves \(\Lambda(Dw_1)\) by its conductor. Gauss decomposition of
\(K_1\) gives these three factors, since all pivots
are units. Averaging over \(K_1\) and applying the
\(\chi_t\)-projector preserve the nonzero evaluation.

To prove uniqueness, let \(w\in W^{K_1}[\chi_t]\)
have \(\Lambda(w)=0\) and suppose \(w\ne0\).
Separation by the mirabolic translates gives a translate testing this nonzero vector, even though its original value is zero:
\[
 \Lambda(w)=0,\qquad \Lambda(gw)\ne0\quad\text{for some mirabolic }g.
\] Starting above the last row,
suppose the rows below row \(i\) are already upper
unitriangular. For \(j\le i\) and \(a\in\mathcal O\),
the conjugate of \(1+aE_{j,i+1}\) by \(g\) is upper
unipotent. Its only possibly nonzero simple entry
is \(ag_{ij}\) in position \((i,i+1)\), by the
support of \(ge_j\) and of the \((i+1)\)-st row of
\(g^{-1}\). Evaluating \(\Lambda(g x_{j,i+1}(a)w)\) in these two ways and dividing by
\(\Lambda(gw)\ne0\) gives
\[
 \psi(\widetilde t_i a g_{ij})
     =\psi(\widetilde t_i a\delta_{ij})\quad(a\in\mathcal O).
\]

As \(\widetilde t_i\) is a unit, the conductor condition gives
\(g_{ij}-\delta_{ij}\in\mathfrak p\).
Right multiplication by \(1+e_i b^{\mathsf T}\),
with \(b_j=-g_{ij}/g_{ii}\) for \(j<i\),
\(b_i=g_{ii}^{-1}-1\) and \(b_j=0\) for \(j>i\),
is a correction in \(K_1\). It makes row \(i\)
upper unitriangular and preserves the lower rows
and the nonzero evaluation. Let \(k_1\in K_1\) be the product of the successive right corrections.
The final matrix is \(u=gk_1\in U\). Since \(k_1 w=w\), this gives
\[
 0\ne\Lambda(gw)=\Lambda(gk_1w)=\Lambda(uw)
                   =\psi_U(u)\Lambda(w)=0,
\]
where \(\psi_U\) is the chosen Whittaker character on \(U\).
Thus \(\Lambda\) is injective on each nonzero space
\(W^{K_1}[\chi_t]\); each has dimension one.
Summing these dimensions by \eqref{eq:depth-one-projector-splitting} gives
\[
                         D_0=(q-1)^{n-1}.
\]

\proofstep{The fixed tame coefficient type.}
At a split tame place, the fixed coefficient type
of Section~\ref{sec:count-two-level-data} gives
zero monodromy and fixed inertia labels. These are
the same Teichm\"uller lifts for both points.
The labels extend to the Weil group, so Frobenius
preserves every inertia isotypic space. Frobenius
semisimplification expresses the parameter as a
sum of Weil characters with exactly those inertia
labels. The local components of cuspidal transfers
are generic.

By the segment classification,
zero monodromy gives singleton segments, and the
one-dimensional Weil characters give rank-one
cuspidal factors \cite[Theorem 9.7]{zelevinsky}.
Thus the local component is the irreducible
normalized principal series of the resulting
characters. For \(K_w=\mathrm{GL}_n(\mathcal O_{L_w})\) and upper
triangular subgroup \(B_w\), Iwasawa decomposition gives
\[
 \left.\operatorname{Ind}_{B_w}^{\mathrm{GL}_n(L_w)}
                      (\textstyle\bigotimes_i\chi_i)\right|_{K_w}
   \simeq
 \operatorname{Ind}_{B_w\cap K_w}^{K_w}
                         (\textstyle\bigotimes_i\chi_i|_{\mathcal O_{L_w}^\times}).
\]
The modulus factor is trivial on \(B_w\cap K_w\).
The unit characters agree by the stipulated inertia
and local reciprocity. Tensor the compact isomorphisms
and take invariants with the same fixed tame coefficient
representation. This gives the common \(d_{\mathrm{tame}}>0\).

\proofstep{The ordered uniformizer values of the ordinary line.}
At a place over \(\ell\), put \(q_\ell=\ell^f\),
where \(f\) is the residue degree. For the
unnormalized translations by
\(t_j=\operatorname{diag}(\varpi I_j,I_{n-j})\)
write \(\mu_j\) for the prescribed eigenvalues and
\(\mu_0=1\). Write
\(a_i=\prod_\tau \tau(\varpi)^{\lambda_{\tau,n-1-i}}\)
using the dominant weights with indices from zero.
Under the isomorphism from the \(\ell\)-adic
coefficient field to \(\mathbb C\), \(\mu_j\)
corresponds to \(\prod_{i<j}a_i\) times a unit. Dividing
successive products gives
\[
  \alpha_i=q_\ell^i\mu_{i+1}/\mu_i,\qquad
  v_\ell(\alpha_i)=if+v_\ell(a_i).
\]
Dominance makes \(v_\ell(a_i)\) nondecreasing, so the
valuations of the \(\alpha_i\) are strictly increasing.

\proofstep{Injectivity of the projection to the Jacquet module.}
For completeness, write \(U_0=U(\mathcal O)\),
and act on \(U_0\)-invariants by the sum of the
translations \(ut_j\) over
\(U_0/t_jU_0t_j^{-1}\). The product of these operators
is the translation sum for \(t=t_1\cdots t_n\), by
using representatives for the successive
subgroups. If \(\mathcal U_j\) is the translation sum for \(t_j\),
the product \(\prod_j\mathcal U_j\) is invertible on
the joint generalized eigenspace because each eigenvalue
\(\mu_j\) is nonzero. On the other hand any vector of
zero Jacquet image can be written
\[
                      v=\sum_m (u_m-1)w_m .
\]
Choose \(r\) with every \(u_m\in t^{-r}U_0t^r\).
Averaging over this subgroup kills each summand and hence
kills \(v\). After translation
by \(t^r\), this says that the corresponding
translation sum is zero. Consequently the
Jacquet map is injective on the refined
generalized eigenspace.

\proofstep{Uniqueness of the selected torus ordering.}
The exceptional projector and tame type have left the order of the torus characters as the remaining possible multiplicity. The ordinary eigenvalues distinguish that order.

On the normalized Jacquet module the scalar multiplying the torus action of \(t_j\) is
\[
 q_\ell^{j(n-j)}\delta_B(t_j)^{1/2}=q_\ell^{j(n-j)/2}.
\]
If the Jacquet module of an irreducible
component is nonzero, it has a character
quotient. Frobenius reciprocity embeds the
component into the corresponding principal
series. Exactness of the Jacquet functor and
the geometric lemma bound the dimension of
any refined eigenspace by the number of
permutations of the inducing characters with
the prescribed torus eigenvalues
\cite[\S2.12]{bernstein-zelevinsky}.
Write each inducing character as
\(\psi_i=\omega_i|\cdot|^{(n-1)/2}\).

For a permutation \(\sigma\), inserting the modulus into the normalized operator gives the line eigenvalue
\[
 q_\ell^{-j(j-1)/2}\prod_{i<j}\omega_{\sigma(i)}(\varpi),\qquad \delta_B(t_j)=q_\ell^{-j(n-j)} .
\]
Thus the first \(j\) selected \(\omega\)-values
at the uniformizer must have product
\(q_\ell^{j(j-1)/2}\mu_j\). Successive quotients
give the \(\alpha_i\), which are pairwise
distinct. There is at most one such
permutation, even before testing its
characters on the units. This gives the
claimed upper bound of one. Positivity on
the initial contributing constituent
forces equality there and throughout its
split component comparison.

\end{proof}

\subsection{The same projected Hecke module at both points}

The local dimensions can now be multiplied with the global packet multiplicity. Comparing two specializations requires the same projector on the same integral Hecke module, since a change of projector could change that product.

\begin{construction}
\[
\begin{array}{lr}
\text{contribution}&\text{factor in }\mu_e(x)\\ \hline
\text{global multiplicity at the fixed weight}&a_0(\lambda^\infty)\\
\text{ordered eigenspace above }\ell&1\ \text{at each place}\\
\text{tame type}&d_{\mathrm{tame}}\\
\text{exceptional depth-one projector}&D_0=(q-1)^{n-1}
\end{array}
\]
\caption{Factors of the projected multiplicity at an irreducible arithmetic point. At fixed weight and local data each contributes the same factor at both points being compared.}\label{fig:ordinary-multiplicity-factors}
\end{construction}

\begin{proposition}[Multiplicity at fixed algebraic weight]\label{prop:ordinary-multiplicity-equal}
For an irreducible arithmetic point \(x\) of the localized
ordinary support, the selected multiplicity is
\[
 \mu_e(x)=a_0(\lambda^\infty)d_{\mathrm{tame}}D_0>0 .
\]
The local type, exceptional projector, level and algebraic
weight exponents determine these factors. In particular the
multiplicity is the same at any other such point \(y\) with
those data, even if its finite-order characters differ.
\end{proposition}

\begin{proof}
\proofstep{Compatibility of control with the projector.}
Let \(M\) be the finite free localized ordinary
module over the weight ring \(\mathcal W\),
at the level used for the determinant quotient.
The determinant ring maps onto the Hecke
image acting on \(M\). The kernel is the
annihilator of \(M\). Since the projector commutes with both actions, the same splitting is defined over \(\mathcal W\) and over the determinant ring:
\[
 M=eM\oplus(1-e)M .
\]

For an arithmetic point \(x\) whose kernel contains
the Hecke ideal, specialize \(M\) by the restriction
of \(x\) to \(\mathcal W\). The arithmetic
control map is Hecke-equivariant and gives an
identification with the specified classical
ordinary space after extension of scalars.
The classical joint generalized eigenspace
is \(Z_x\), and it is nonzero for \(x\)
in the Hecke support. Write \(C\) for the Hecke-equivariant control map. Its compatibility with the integral translation sum gives
\[
 C\,(|U(k)|e)=(|U(k)|e)\,C .
\] Dividing by the unit
\(|U(k)|\) gives the equality of selected
multiplicities before and after control.
The spherical operators and the ordinary
operators generate the Hecke action. Testing
their generalized eigenvalues therefore gives
the specialized local Hecke summand, and
projecting commutes with this test.

\proofstep{Multiplication of the three local dimensions.}
The same idempotent now acts at the two specializations. The three local dimensions can therefore be multiplied with the global packet count to compare the two projected multiplicities.

Now each contributing index has projected
dimension \(d_{\mathrm{tame}}D_0\), which is
positive in view of its unprojected
nonvanishing. Proposition~\ref{prop:ordinary-stable-constituent-count}
now gives
\[
             \mu_e(x)=a_0(\lambda^\infty)
                             d_{\mathrm{tame}}D_0>0.
\]
At a second irreducible arithmetic point
\(y\) in the same Hecke support with the same algebraic
weight exponents, the stable transfer may change. The
finite-order characters can change as well. Ordinary
control uses the homomorphism from the weight ring
at each point separately, on an open subgroup where
its finite-order factors are trivial.
All three factors on the right agree.
In particular \(\mu_e(x)=\mu_e(y)\) on
the idempotent range which controls the
fixed-weight derivations. The comparison
requires a common algebraic weight, level, tame type
and exceptional projector. The constant
\(a_0\) is fixed by the field, rank and
archimedean weight. The proof allows the
depth of classical specialization and
the auxiliary trace tests to vary.
\end{proof}
